\section{The normal operator of a maximal surface representation}\label{sec:cayley}
Let \(C\) be a closed Riemann surface of genus \(g\ge2\). Let a flat symplectic four-dimensional fiber germ have a marked horizontal section over \(C\), and give its flat horizontal and vertical distributions the split almost-complex structures. Assume the normal symplectic bundle has first Chern number \(2g-2\). Write \(P_\rho(Q)\) for the series of compact open insertions of the multiples of the section in the smooth universal cycle algebra \(S\), with constant term one. Here \(\rho\) denotes the derivative holonomy. A nonlinear germ deforms to this linear model without losing isolation in any degree. This follows from uniform escape of nonzero transverse orbits and exactness of the vertical form on a tube.

Fix a theta characteristic \(\theta^2=K_C\). In the harmonic reduction of \(\rho\), write \(V\) for its positive-degree holomorphic rank-two bundle and \(\gamma:V\longrightarrow K_C V^*\) for its symmetric lower Higgs block. Maximality makes \(\gamma\) an isomorphism. The resulting orthogonal rank-two bundle \(F=V\otimes\theta^{-1}\), with \(\gamma\) untwisted, is the orthogonal bundle underlying its Cayley partner. Write \(w=w_1(F)\) and \(e=\langle w_2(F),[C]\rangle\). If \(w=0\), the two isotropic lines split \(F\) holomorphically as \(N\oplus N^{-1}\). Put \(c=|\deg N|\). The reduction, the component classification and these characteristic classes are those of Bradlow--García-Prada--Gothen; the results here concern the resulting classes of local curve counts \cite{bgg}.

\emph{Theorem.} The connected genus-\(g\) degree-one test of \(P_\rho\) is \((-1)^{h^0(C,V)}\). On the components above it equals
\[\begin{cases}
 (-1)^c,&w=0,\\
 (-1)^{q_\theta(w)+e},&w\ne0,
 \end{cases}
 \qquad
 q_\theta(a)=h^0(\theta\otimes\ell_a)-h^0(\theta)\pmod2 ,\]
where \(\ell_a\) is the flat sign line of \(a\in H^1(C;\mathbb F_2)\). When \(w=0\), the full series in \(S\) equals the image of the holomorphic antidiagonal local series for a pair of line bundles of degrees \(g-1+c\) and \(g-1-c\) with product \(K_C\). When \(w\ne0\), its pullback \emph{as a neighborhood} to its connected orientation double cover has full series \(E_{2g-1}\), the balanced elementary factor there. More generally the genus-\(2g-1\) degree-one test on the connected double cover classified by \(0\ne a\in H^1(C;\mathbb F_2)\) is \((-1)^{\langle a\smile w,[C]\rangle}\). These last signs recover \(w\). For \(w\ne0\) the downstairs sign then recovers \(e\) once \(\theta\) is fixed.

The expression of \(h^0(\theta\otimes F)\) modulo two in terms of \(w_1,w_2\) is the rank-two case of the established orthogonal spin index formula. The geometric identification here is with the determinant of the suspended core and its universal degree-one test \cite{bspectral}.

The sign gives actual distinctions among series in \(S\). For example, two copies of one uniformizing spin lift give positive degree-one sign. The direct sum of lifts belonging to opposite spin parities has negative degree-one sign. At its least source genus equality of the two degree-one insertions would identify two regular marks with opposite determinant orientations. The multiplicity-one detection theorem \cite[Theorem~5.1]{WueProduct} excludes this equality. No identification of the full negative-sign series with the inverse of the positive one follows from that test \cite{WueProduct}.

\subsection{The real normal operator from the harmonic reduction}

Every maximal representation here is completely reducible. Indeed the elementary Euler bound excludes any nonzero invariant isotropic subspace; the radical of any invariant subspace would be such a subspace. The orthogonal complement therefore splits off each invariant subspace. Nonabelian Hodge theory supplies a harmonic reduction to \(\mathrm U(2)\). Its compatible vertical complex structure may replace the original one along a path of compatible linear structures: the flat fiber form still makes all compact curves in the tube horizontal. The associated rank-two bundle \(V\) has \(\deg V=c_1(W)[C]=2g-2\). In a unitary frame the complexification of the real standard representation is \(V\oplus\bar V\), where the Hermitian metric identifies \(\bar V\) with \(V^*\). A real vector with first coordinate \(s\) has second coordinate \(\bar s\). In these coordinates the complex flat connection is \(d_A+\Phi+\Phi^*\) with \(A\) unitary and
\[\Phi=\begin{pmatrix}0&\beta\\ \gamma&0\end{pmatrix}.\]
In particular the \((0,1)\) operator on the first coordinate of a real section is
\[D s=\bar\partial_V s+R_h(s),\qquad
 R_h(s)=h^\sharp\bigl(\gamma(s,\cdot)\bigr).\]
Here \(h^\sharp:V^*\to V\) is the conjugate-linear metric identification; it conjugates the differential on the canonical factor as well. To check this independently in a flat real connection, its Cartan decomposition in unitary frames has a unitary part \(A\) and a symmetric complex-antilinear part \(B\circ\mathrm{conj}\). Write the first-row coefficient of the latter as \(P\,dz+R\,d\bar z\). The conjugate row has \((1,0)\) term \(\bar R\,dz\), namely \(\gamma\). Symmetry gives \(R=\overline{\gamma_z}\) in this unitary frame. The \((0,1)\) projection on the vertical complex coordinate is exactly \(\bar\partial_A s+R\bar s\,d\bar z\). This avoids a dependence on a choice of scalar in the identification of the standard symplectic form with its block matrix. The Hermitian sign and \(V\) convention can be checked on a uniformizing plane, where it gives the degree-\(g-1\) spin line \cite{bgg}.

The graph equation for a smooth section of the flat bundle is \(Ds=0\). More generally a graph over a holomorphic map \(f:D'\to C\) satisfies the pullback equation (in the zeroth-order term the base differential is pulled back). Conversely the holomorphic structure on \(V\) and any smooth conjugate-linear \(R:V\to\Lambda^{0,1}C\otimes V\) determine a fiberwise-linear almost-complex structure on its total space with holomorphic projection and standard complex fibers. In bundle coordinates its \((0,1)\) directions are the vertical \((0,1)\) directions and a lift of \(\partial_{\bar z}\) whose vertical \((1,0)\) part is \(-R_{\bar z}(s)\). This is intrinsic because a change of holomorphic bundle frame gives the Dolbeault transformation rule. Graphs satisfy \(\bar\partial s+R(s)=0\). The harmonic flat structure corresponds to the displayed \(R_h\). These structures are tamed on a common small disk bundle in any compact parameter family: use a closed form consisting of a sufficiently positive base area form and a fiberwise symplectic form obtained by differentiating the metric radial primitive with a unitary connection. On the zero section the horizontal-vertical splitting is complex; tameness persists by compactness on a tube. Its first Chern number stays zero on the section and hence throughout the tube.

For any Hermitian metric \(h\) and \(t>0\) the structure with equation \(D_t=\bar\partial_V+tR_h\) has only the compact reduced core. Indeed on the normalization of a compact curve projecting nonconstantly by \(f\), its fiber coordinate is a section \(s\in\Gamma(f^*V)\). The symmetric form paired with \(s,s\) and then pulled through \(df\) gives a globally defined \((1,0)\)-form. In a local coordinate \(u\) on the normalization, with \(\gamma=G\,dz\) in unitary frames, its \(\bar\partial\) is
\[-2t\,\|G(f)s\|^2\, |f_u|^2\,d\bar u\wedge du .\]
To compute in unitary frames use the metric compatibility of \(A\) and \(\bar\partial_A\gamma=0\); equivalently differentiate the holomorphic symmetric form and contract both copies of \(\bar\partial_{f^*V}s\). In nonunitary holomorphic frames the same expression includes the positive inverse Hermitian metric in the norm. Stokes' theorem forces \(s\) to vanish on the unramified open set and then everywhere, since \(G\) is invertible. A compact curve in a single complex affine fiber is constant. In particular the only degree-\(m\) cycle is \(m[C]\) for all positive \(t\), including through an arbitrary smooth homotopy of the metric. This argument uses actual cycles and every branched normalization.

\subsection{Determinant signs}

For the identity map \(C\to C\), the complex-linear operator \(D_0=\bar\partial_V\) has index zero. Its kernel is \(H^0(V)\) and its complex cokernel is \(H^1(V)\). By the isomorphism \(\gamma\) and Serre duality the cokernel is the dual of \(H^0(V)\). The first-order map from kernel to cokernel for \(D_t\) is induced by \(R_h\). Pairing its value at \(s\) with the dual section \(\gamma(r,\cdot)\) gives, up to a common nonzero scalar from the area convention, the integral of the Hermitian inner product of \(\gamma(r,\cdot)\) and \(\gamma(s,\cdot)\). This is a nondegenerate conjugate-linear identification. Splitting kernel and cokernel off by complex complements gives a finite-dimensional reduced operator \(t[R_h]+O(t^2)\); the complementary block remains invertible with its complex orientation. The real determinant of a conjugate-linear isomorphism of complex \(d\)-planes has sign \((-1)^d\). The Stokes argument makes \(D_t\) invertible for every \(t>0\), so the computed sign persists at \(t=1\). This is the determinant sign of the isolated regular graph, since varying the source cancels the tangent operator of the degree-one map to \(C\). Its automorphism group is trivial and the source of genus \(g\) leaves no contracted handles or nodes. The construction generalizes the Lee--Parker real spectral-flow computation for one spin line \cite{lp} \cite{dwcount}.

If \(w=0\) the oriented orthogonal group in complex rank two is \(\mathbb C^\times\) acting in reciprocal characters. Hence \(V=\theta N\oplus\theta N^{-1}\). By Serre duality the two zeroth cohomology dimensions are \(h^0(\theta N)\) and \(h^1(\theta N)\). Their sum modulo two is \(\chi(\theta N)=\deg N\). This computes the signs on the balanced component, the exceptional oriented components \(1\le c\le2g-3\), and the Hitchin components \(c=2g-2\), with the standard integer convention of the classification. In particular all Hitchin components have positive degree-one test; individual odd \(c\) have negative test.

For \(w\ne0\) use a split orthogonal representative \(\ell_b\oplus\ell_{w+b}\) with \(\langle b\smile w,[C]\rangle=e\). One exists for either \(e\) by nondegeneracy of the surface cup product. These examples are the tensor products of a fixed uniformizing lift for \(\theta\) with two real sign characters. Their Cayley invariants are \(w_1=w\) and \(w_2=b\smile(w+b)\) by the Whitney formula and \(b^2=0\). They represent the components labeled by those invariants. Alternatively within flat \(\mathrm O(2)\) twists, after fixing the determinant, representations are circle cocycles with action by this determinant. The exponential sequence of twisted coefficients shows that the obstruction distinguishing components is the twisted Euler class in \(H^2(C;\mathbb Z_w)=\mathbb Z/2\); equal classes differ by a cocycle lifting to a real cocycle and can be joined by scaling the lift, up to a rotation coboundary. Its mod-two reduction is \(w_2\) \cite{bgg}.

On this split model \(h^0(V)=h^0(\theta\ell_b)+h^0(\theta\ell_{w+b})\). The Riemann--Mumford identity \(q_\theta(b+w)=q_\theta(b)+q_\theta(w)+\langle b\smile w,[C]\rangle\) proves the sign formula. Negative and positive signs occur even at fixed nonzero \(w\). The sign persists on a representation component by the geometric continuation above, or directly by continuity of the invertible normal operator. The same calculation gives the stated two-spin test without invoking the component classification \cite{mumford}.

\subsection{All multiples for the oriented Cayley bundle}

Write \(L_+=\theta N\), \(L_-=\theta N^{-1}\), so \(L_+L_-=K_C\). In this splitting \(\gamma\) pairs the two lines and vanishes on each individual line. Move the harmonic metric through Hermitian metrics to one making the lines orthogonal, using \(\bar\partial_V+R_h\) throughout. Stokes keeps every multiple isolated; at this stage no circle equivariance is needed. For the final metric \(R_h\) sends the conjugate of each line to the other. Hence the circle acting on \((L_+,L_-)\) with weights \((+1,-1)\) commutes with \(D_t=\bar\partial_V+tR_h\) for every \(0\le t\le1\). The endpoint at \(t=0\) is the ordinary holomorphic total space of \(L_+\oplus L_-\). At every \(t>0\) there are only the core cycles.

For completeness fix \(m\) and choose enough base evaluations of the elementary symmetric polynomials of the two coordinate multisets that their common zero in the integrable cycle space is \(m[C]\). Each order-\(i\) coefficient is a holomorphic section of \(L_+^i\) or \(L_-^i\), by the characteristic polynomial on the finite direct image of the normalization with cycle multiplicities. More than \(m\max(0,\deg L_+,\deg L_-)\) distinct points suffice. The resulting test weights are \(+i\) and \(-i\). Let \(n\) be the total number of scalar tests. On the finite circle Borel stage \(\mathbb P^n\), take their Thom product divided by the product of their signed weights. It defines the complex antidiagonal degree-\(m\) insertion for these two line bundles. The original balanced construction carries over with this explicit degree bound. It depends only on the unordered degrees and genus, by continuation in connected Picard varieties and in the moduli of smooth genus-\(g\) curves, using sufficiently many moving evaluations locally and tuple stabilization on overlaps. One may exchange the labels by reindexing and reversing the circle generator; in the Borel quotient complex conjugation on \(\mathbb P^n\) has degree \((-1)^n\), which cancels reversal of the \(n\) signed weights. Equivalently fix the larger degree as the positive label throughout \cite{WueBalanced}.

In a small invariant disk over the whole interval the relative cycle space is the union of its integrable endpoint and the interval of core cycles. Both are closed subspaces by the compact-cycle topology. The map which fixes that endpoint and collapses the core interval onto \(m[C]\) is continuous by pasting and equivariant. Its inverse image of a compact subset is compact, being the union of that subset and a compact interval if the core belongs to it. On the finite Borel stage the analogous proper retraction pulls back the normalized Thom insertion with compact support. The invariant metric gives the associated disk bundle; its integrable endpoint may equally be identified, by open excision at the compact supported core locus, with the holomorphic Borel family of the full split bundle. On \(t>0\) the restriction is the ordinary Euler class divided by the signed weights, namely \(h^n\) with \(h=c_1(\mathcal O_{\mathbb P^n}(1))\). Its pushforward is the unit of the isolated multiple at a point. The smooth-family interval relation therefore identifies each original coefficient with the image of the corresponding complex antidiagonal coefficient. The homotopy from the original harmonic metric preceded this equivariant interval, so its failure to preserve the circle is immaterial \cite{WueProduct}.

In particular \(c=0\) gives \(E_g\) as in Section~\ref{sec:horizontal}. For \(c>0\) the result is a geometric comparison with the \emph{unbalanced} antidiagonal construction. Its algebraic numerical degree-two residue is the balanced one times \(\cos(cu)\) by the Bryan--Pandharipande local-curve formula. Distinguishing the full smooth classes through a numerical evaluation requires a justified comparison of that evaluation in its stated domain. The geometric identification just proved uses no such evaluation \cite{WueBalanced}.

\subsection{Covers and the limits of the comparison}

For \(w\ne0\), pullback to the connected orientation double cover \(p:\widehat C\to C\) splits \(p^*F=N\oplus N^{-1}\). Its deck involution exchanges the two isotropic lines, so \(\deg N=-\deg N\) and \(\deg N=0\). The pullback theta characteristic squares to \(K_{\widehat C}\) since \(p\) is étale. The pulled-back harmonic reduction solves the harmonic equations, and the preceding proof gives the balanced factor \(E_{2g-1}\). This statement concerns the neighborhood on the covering base. A circle rotating the two ordered lines is sent to its inverse by the deck involution. The argument does not descend that circle construction or determine the full downstairs series on the components with \(w\ne0\).

For an arbitrary connected double cover \(p_a\) the normal determinant sign can be computed at the split sign representative when \(w\ne0\). Since \(p_{a*}\mathcal O=\mathcal O\oplus\ell_a\), its parity exponent is the sum of the four parities for spin lines indexed by \(b,w+b,b+a,w+b+a\). The quadratic identity shows that this sum is \(\langle a\smile w,[C]\rangle\). If \(w=0\), the oriented split integer doubles under pullback and again gives the asserted sign. Thus the full collection of connected double-cover degree-one tests recovers \(w\) by nondegeneracy of the intersection form. Within the nonoriented components these tests together with the downstairs degree-one test recover the two component labels relative to \(\theta\). These are tests on a system of neighborhoods under base change. They make no assertion that a single downstairs series distinguishes every component.

\emph{Corollary for the unmarked base.} Among the components with \(w\ne0\), the unmarked downstairs series has precisely two values. Its degree-one sign determines which of the two occurs. Indeed, for each fixed nonzero \(w\) the component classification gives two components, and the two signs distinguish them. An orientation-preserving diffeomorphism pulls the flat target and its section isomorphically to the corresponding neighborhood over the pulled-back base. Changing that complex structure back along a path preserves the full unmarked series: the exact vertical form and isolation argument are unchanged, and the tangent first Chern number is constant. Such diffeomorphisms act transitively on nonzero mod-two cohomology classes of a closed oriented surface (symplectic changes of homology basis realize this transitivity). They preserve the sign of the regular core. Consequently all the components of a given sign are equivalent for this unmarked series, and opposite signs give different values. This does not give the higher coefficients of either value. The earlier detection of \(w\) uses the indexed collection of covers of a marked base.
